\begin{textsample}{POS Dim 3 – pb – Score 23.00 – pb/pb\_inf\_9.txt}  \label{ex:f3_pos_085}
3.5.1.3. \textbf{ESCUTA}, FALA, PENSAMENTO E IMAGINAÇÃO Sabemos que as \textbf{crianças} aprendem a se comunicar na experiência com \textbf{parceiros}, através de diálogos, imitação, \textbf{gestos}, o que lhes favorece a construção de sentidos.
Na etapa que corresponde à Educação \textbf{Infantil}, as \textbf{crianças} estão em pleno desenvolvimento da oralidade, bem como do uso de diversas formas de comunicação, sendo as \textbf{instituições} de Educação \textbf{Infantil} também ambientes sociais onde aprendem a falar, a se expressar, a exprimir -se por meio de diferentes linguagens.
Portanto, precisam ser oportunizadas significativas situações que lhes permitam a \textbf{escuta}, a fala, a comunicação de um modo geral, que envolvam o pensamento e a imaginação, a partir de experiências que considerem seus olhares, \textbf{gestos}, choros, suas conversas, narrativas construídas individualmente ou em grupo, ampliando assim seu vocabulário e favorecendo a apropriação da sua língua materna.
Nós, professores/as, devemos mediar o desenvolvimento da comunicação da \textbf{criança}, na organização da \textbf{rotina} e na elaboração do planejamento favorecendo situações que lhe possibilitem o interesse pelas imagens, pelas cantigas, \textbf{brincadeiras} de \textbf{roda}, jogos cantados e outros textos do universo popular e clássico.
Ao lermos para elas, oportunizando que construam hipóteses sobre a narrativa, fazendo antecipações e/ou suposições sobre o que vai acontecer, imaginando ou descrevendo cenas, realizem contagem, nomeiem objetos e relatem ações do cotidiano, estamos atuando neste campo.
As \textbf{crianças} devem vivenciar desde cedo situações em que podem usar a linguagem de forma criativa, para expressar ideias e emoções e interagir com os demais.
A experiência contínua da \textbf{escuta} de histórias permitirá que as \textbf{crianças} desenvolvam comportamento leitor, que envolve desde o prazer com o \textbf{som} da narrativa até a construção de uma concepção sobre a língua escrita.
Nesse processo, a \textbf{imersão} se dá a partir do próprio universo da \textbf{criança}, daquilo que ela já traz como vivência em seu grupo familiar e cultural e se amplia com a \textbf{curiosidade} e a vivência de situações que as levem a refletir, formulando hipóteses sobre a escrita, dialogando com \textbf{crianças} e com \textbf{adultos}, participando de situações de leitura e escrita, aprendendo a ler o mundo das imagens, das letras, dos números, das palavras e dos textos.
Essas experiências devem ser organizadas de forma que as \textbf{crianças} observem situações reais de escrita e seu uso social, como também sejam oportunizadas a escrita como expressão espontânea.
Assim, a leitura e a escrita não devem ser exploradas mecanicamente.
A leitura deve ser uma fonte de prazer, de desenvolvimento da imaginação e de aprendizagem e as práticas de oralidade e escrita devem ocorrer de forma contextualizada e significativa, em um ambiente que favoreça o convívio com o universo da imaginação, da expressão oral e escrita de ideias e sentimentos e do próprio movimento de \textbf{escuta} dos outros com quais a \textbf{criança} convive.
Desde o \textbf{nascimento}, as \textbf{crianças} participam de situações comunicativas cotidianas com as pessoas com as quais interagem.
As primeiras formas de interação do \textbf{bebê} são os movimentos do seu corpo, o olhar, a postura corporal, o sorriso, o choro e outros recursos vocais, que ganham sentido com a interpretação do outro.
Progressivamente, as \textbf{crianças} vão ampliando e enriquecendo seu vocabulário e demais recursos de expressão e de compreensão, apropriando -se da língua materna – que se torna, pouco a pouco, seu veículo privilegiado de interação.
Na Educação \textbf{Infantil}, é importante promover experiências nas quais as \textbf{crianças} possam falar e \textbf{ouvir}, potencializando sua participação na cultura oral, pois é na \textbf{escuta} de histórias, na participação em conversas, nas descrições, nas narrativas elaboradas individualmente ou em grupo e nas implicações com as múltiplas linguagens que a \textbf{criança} se constitui ativamente como sujeito singular e pertencente a um grupo social.
Desde cedo, a \textbf{criança} manifesta \textbf{curiosidade} com relação à cultura escrita : ao \textbf{ouvir} e acompanhar a leitura de textos, ao observar os muitos textos que circulam no contexto familiar, comunitário e escolar, ela vai construindo sua concepção de língua escrita, reconhecendo diferentes usos sociais da escrita, dos gêneros, suportes e portadores.
Na Educação \textbf{Infantil}, a \textbf{imersão} na cultura escrita deve partir do que as \textbf{crianças} conhecem e das \textbf{curiosidades} que deixam transparecer.
As experiências com a literatura \textbf{infantil}, propostas pelo educador, mediador entre os textos e as \textbf{crianças}, contribuem para o desenvolvimento do gosto pela leitura, do \textbf{estímulo} à imaginação e da ampliação do conhecimento de mundo.
Além disso, o contato com histórias, contos, fábulas, poemas, cordéis etc. propicia a familiaridade com \textbf{livros}, com diferentes gêneros literários, a diferenciação entre ilustrações e escrita, a aprendizagem da direção da escrita e as formas corretas de manipulação de \textbf{livros}.
Nesse convívio com textos escritos, as \textbf{crianças} vão construindo hipóteses sobre a escrita que se revelam, inicialmente, em rabiscos e garatujas e, à medida que vão conhecendo letras, em escritas espontâneas, não convencionais, mas já \textbf{indicativas} da compreensão da escrita como sistema de representação da língua.
Venho de uma família em que se \textbf{contava} muita história.
Com \textbf{livro} ou sem \textbf{livro}.
E os repertórios variavam muito, de acordo com o contador.
Minha mãe era especialista em contos de fada.
Meu pai sempre trazia uns clássicos diferentes, muitas vezes mostrando as figuras nuns livrões que tirava da estante.
Minha avó \textbf{contava} as histórias populares de nossa tradição oral, cheias de almas do outro mundo, heróis bobos ou espertalhões, bichos que falavam, gigantes...
Entre elas, talvez as minhas preferidas fossem as de Pedro Malasartes, de que ela parecia ter um estoque interminável.
As \textbf{crianças} vivem inseridas em espaços e tempos de diferentes dimensões, em um mundo constituído de fenômenos naturais e socioculturais.
Desde muito \textbf{pequenas}, elas procuram se situar em diversos espaços ( rua, bairro, cidade etc. ) e tempos ( \textbf{dia} e noite ; hoje, ontem e amanhã etc. ). \textbf{Demonstra} também \textbf{curiosidade} sobre o mundo físico ( seu próprio corpo, os fenômenos atmosféricos, os animais, as plantas, as transformações da natureza, os diferentes tipos de materiais e as possibilidades de sua manipulação etc. ) o mundo sociocultural ( as relações de parentesco e sociais entre as pessoas que conhece ; como vivem e em que trabalham essas pessoas ; quais suas tradições e seus costumes ; a diversidade entre elas etc. ).

% matched lemmas: adulto, bebê, brincadeira, contar, criança, curiosidade, demonstrar, dia, escuta, estímulo, gesto, imersão, indicativo, infantil, instituição, livro, nascimento, ouvir, parceiro, pequeno, roda, rotina, som
\end{textsample}
